% (C) 2026 Ralf Stephan, in collaboration with Claude Code.  CC0 1.0.
\chapter{Why carry graphs}\label{ch:overview}

\section{One recurrence}

Let \(p>q\ge 1\) be coprime integers, \(\xi>0\) a real number and \(\nu\) a real shift.
Write
\[
  \xi\Bigl(\frac pq\Bigr)^{n}+\nu \;=\; x_n+y_n,\qquad x_n\in\Z,\quad y_n\in[0,1),
\]
so that \(x_n=\lfloor \xi(p/q)^n+\nu\rfloor\) is the integer part and \(y_n\) the fractional
part. Multiplying the identity at \(n\) by \(p/q\) and comparing with the identity at \(n+1\)
gives
\begin{equation}\label{eq:carry}
  q\,x_{n+1} \;=\; p\,x_n + s_n,
  \qquad
  s_n \;=\; p\,y_n - q\,y_{n+1} - (p-q)\,\nu \in \Z .
\end{equation}
The integer \(s_n\) is the \emph{carry} of step \(n\). It takes only finitely many values,
because \(y_n,y_{n+1}\in[0,1)\), and it is the only information that passes from the
fractional part to the integer part. For \(q=1\) and \(\nu=0\) the recurrence is
\(x_{n+1}=p\,x_n+s_n\) with \(s_n\in\{0,\dots,p-1\}\), and \(s_n\) is the \(n\)-th base-\(p\)
digit of \(\xi\) after the point.

Many old questions about the integers \(\lfloor\xi(p/q)^n\rfloor\) and the fractional parts
\(\frc{\xi(p/q)^n}\) concern the \emph{carry word} \(s_0s_1s_2\cdots\):
\begin{itemize}
\item whether \(x_n\) is composite for infinitely many \(n\) \cite{FS67,DN05,Dub09};
\item whether every \(x_n\) can avoid a finite set \(\mathcal P\) of prime divisors;
\item whether all \(y_n\) can stay in a short interval
  (Mahler's \(Z\)-number problem \cite{Mah68,FLP95}).
\end{itemize}
In each case the question can be asked about finite data: a residue of \(x_n\) modulo some
\(M\), the cell of \(y_n\) in a partition of \([0,1)\) into \(D\) intervals, or both. The
recurrence~\eqref{eq:carry} moves this finite data along a finite labelled digraph whose arc
labels are the carries. This book is about that digraph.

\section{The object}

\begin{definition}
A \emph{carry graph} over a label set \(A\) is a set \(V\) of vertices together with a map
\(\succc\colon V\times A\to\{\text{finite subsets of }V\}\). There is an arc
\(v\xrightarrow{a}w\) when \(w\in\succc(v,a)\).
\end{definition}

The vertex set may be infinite. Only the out-neighbourhoods have to be finite, and that is
the condition the rest of the book uses. Forgetting the labels gives an ordinary digraph in the
sense of \cite[§1.2]{BJG09}. Keeping them gives a labelled graph in the sense of
\cite[§3.1]{LM95}. A finite carry graph is therefore an object of symbolic dynamics, and most
of the tools of \cite{LM95} chapters~1--4 apply to it.

Three families of carry graphs recur.
\begin{description}[style=nextline]
\item[Residue graphs.] Vertices are residues \(r\in\Z/M\). Under the label \(s\), the arc
  \(r\xrightarrow{s}(pr+s)q^{-1}\) is present when \(\gcd(q,M)=1\). Here \(x_n\bmod M\) walks
  along this graph and the label is the carry. Deleting the residues divisible by a prime of
  \(\mathcal P\) leaves the residues an avoiding orbit may visit.
\item[Cell graphs.] Vertices are the cells \([c/D,(c+1)/D)\), \(0\le c<D\). There is an arc
  \(c\xrightarrow{s}c'\) when some \(y\) in cell \(c\) has \((p\,y-(p-q)\nu-s)/q\) in cell
  \(c'\). Then \(y_n\) walks along this graph and the label is again the carry.
\item[Products.] The pair \((x_n\bmod M,\ \text{cell of }y_n)\) walks along the \emph{label
  product} of a residue graph and a cell graph. Its arcs are the pairs of arcs carrying the
  same label (Chapter~\ref{ch:splitting}).
\end{description}

\section{Covers, and the shape of every theorem}

The orbit itself lives in an infinite space: the pairs \((x_n,y_n)\), or the reals \(\xi\).
A carry graph \(X\) is tied to the orbit by a \emph{cover}. This is an assignment
\(\rho\colon V\to\mathcal P(S)\) of a region of state space to every vertex, such that every
step of the dynamics from \(\rho(v)\) under carry \(a\) lands in \(\rho(w)\) for some arc
\(v\xrightarrow{a}w\). A cover makes every orbit a path of \(X\), and the carry word of the
orbit the label of that path. Conversely, \(X\) \emph{realises} the dynamics when every finite
labelled path is the carry word of some orbit segment.

With these two notions, every result about the problems above has the same shape:
\[
  \begin{gathered}
  \text{cover of the dynamics by }X\;+\;\text{a property of the finite graph }X\\
  \Longrightarrow\;\text{a property of every orbit.}
  \end{gathered}
\]
The graph properties that occur are few:
\begin{itemize}
\item every infinite path of \(X\) is eventually periodic (\emph{zero entropy},
  Chapter~\ref{ch:entropy}); combined with the aperiodicity of carry words
  \cite[Lemma~2]{DN05}, an avoiding orbit then cannot exist;
\item every infinite path leaves a vertex set after a bounded number of steps
  (\emph{escape}, Chapter~\ref{ch:core});
\item the number of labelled paths of length \(n\) is at most \(C\lambda^n\), which bounds an
  entropy or a Hausdorff dimension (Chapter~\ref{ch:covers});
\item and, in the other direction, the existence of two closed walks of equal length
  through one vertex (a \emph{horseshoe}). When \(X\) realises the dynamics, this produces
  orbits with positive entropy, and it shows that no proof of the first kind exists for that
  graph.
\end{itemize}
A computer search finds the graph and a finite witness for the property, such as a rank
function, a peeling sequence, a sub-eigenvector or a pair of closed walks. A proof assistant
checks the witness, and the transfer theorems of this book turn the checked property into a
statement about orbits. The book develops the graph theory and the transfer theorems once,
for all carry graphs. The witnesses and searches belong to the individual problems and appear
here only as examples.

\section{What is in the book}

Chapters~\ref{ch:graphs}--\ref{ch:core} cover walks, Kőnig's lemma, reachability, strong
components and their acyclic orderings \cite[§1.4--1.5]{BJG09}. They also introduce ranks
and the peeling of a graph down to its core, the part that carries every infinite path.
Chapter~\ref{ch:entropy} counts paths. It proves the dichotomy between polynomial and
exponential growth, the equivalence of six descriptions of zero entropy, and the
eigenvector bounds that stand in for the Perron--Frobenius theorem \cite[§4.1--4.4,
§5.3]{LM95}. Chapter~\ref{ch:sofic} studies labels: label languages, the subset
construction, follower sets and sofic shifts \cite[ch.~3]{LM95}. Chapter~\ref{ch:splitting}
covers state splitting, higher-block and line graphs, homomorphisms, and label products
\cite[§2.3--2.4]{LM95}, \cite[§4.5]{BJG09}. Chapter~\ref{ch:covers} defines covers and
proves the transfer theorems. Chapter~\ref{ch:affine} specialises to the affine carry graphs
of the recurrence~\eqref{eq:carry}. It proves that the residue--cell graph is a label
product and works through the problems listed above as examples. Chapter~\ref{ch:flows} turns from
avoidance to frequency: circulations, the balance lemma for the visit frequencies of a walk,
and why flow conservation alone cannot prove that \(\frc{(3/2)^n}\) visits every interval with
positive density.

\section{Lean}

Every definition and theorem in this book is formalised in Lean~4 on top of Mathlib
\cite{Mathlib}, in the directory that also holds this book. The two use the same names:
a carry graph is the structure \lean{CarryGraph V A} with one field
\lean{succ : V \(\to\) A \(\to\) Finset V}. Each chapter ends with a table that matches its numbered
results to Lean declarations, and Appendix~\ref{app:dictionary} collects these tables.
Appendix~\ref{app:mathlib} lists what Mathlib already provided and what had to be built.

\section*{Notes}
The recurrence~\eqref{eq:carry} and the carry word are classical. The aperiodicity of the
carry word is \cite[Lemma~2]{DN05}. Residue dynamics for \(q=1\) is the method of
\cite{FS67} and \cite{Dub09}. Cell dynamics and the intervals that confine
\(\frc{\xi(p/q)^n}\) go back to \cite{FLP95}. Finite graph witnesses for composite terms of
\(\lfloor\xi(a/b)^n\rfloor\) appear in \cite{Oda26}.
